\documentclass[../../main.tex]{subfiles}

\begin{document}

\section{Tychonoff 乘积定理}\label{section:ch12-Tychonoff}

对由集合 $\Lambda$ 标记的集族 $\{X_{\lambda}\}_{\lambda\in\Lambda}$, 我们定义笛卡儿积 $\prod_{\lambda\in\Lambda}X_{\lambda}$ 为从指标集 $\Lambda$ 到 $\bigcup_{\lambda\in\Lambda}X_{\lambda}$ 的映射族, 使得每个指标 $\lambda\in\Lambda$ 被映到 $X_{\lambda}$ 的某个成员. 对笛卡儿积的成员 $x$ 与指标 $\lambda\in\Lambda$, 习惯上记 $x(\lambda)$ 为 $x_{\lambda}$ 且称 $x_{\lambda}$ 为 $x$ 的第 $\lambda$ 个分量. 对每个参数 $\lambda_{0}\in\Lambda$, 我们定义 $\lambda_{0}$ 的投影映射 $\pi_{\lambda_{0}}:\prod_{\lambda\in\Lambda}X_{\lambda}\to X_{\lambda_{0}}$ 为
\begin{align*}
\text{对 }x\in\prod_{\lambda\in\Lambda}X_{\lambda},\quad\pi_{\lambda_{0}}(x)=x_{\lambda_{0}}.
\end{align*}
我们已定义了两个度量空间的笛卡儿积上的乘积度量. 可用明显的方式推广到有限个度量空间的笛卡儿积上的度量. 此外, 可数个度量空间的笛卡儿积上存在自然的度量(见习题 16).

拓扑空间的有限族的笛卡儿积上的拓扑有自然的定义. 给定拓扑空间族 $\{(X_{k},\mathcal{T}_{k})\}_{k=1}^{n}$, 乘积族
\begin{align*}
\mathcal{O}_{1}\times\cdots\mathcal{O}_{n}\times\mathcal{O}_{n}
\end{align*}
是 $\prod_{1\leqslant k\leqslant n}X_{k}$ 上的拓扑的基, 其中每个 $\mathcal{O}_{k}$ 属于 $\mathcal{T}_{k}$. 由这些基本集的并组成的笛卡儿积上的拓扑称为 $\prod_{1\leqslant k\leqslant n}X_{k}$ 上的乘积拓扑.

对拓扑空间来说, 新颖的地方在于乘积拓扑可在拓扑空间的任意笛卡儿积 $\prod_{\lambda\in\Lambda}X_{\lambda}$ 上定义. 指标集不再要求是有限的, 甚至不要求是可数的.

\begin{definition}[乘积拓扑]\label{definition:ch12-乘积拓扑}
令 $\{(X_{\lambda},\mathcal{T}_{\lambda})\}_{\lambda\in\Lambda}$ 为由集合 $\Lambda$ 标记的拓扑空间族. 笛卡儿积 $\prod_{\lambda\in\Lambda}X_{\lambda}$ 上的乘积拓扑有形如 $\prod_{\lambda\in\Lambda}\mathcal{O}_{\lambda}$ 的基集, 其中每个 $\mathcal{O}_{\lambda}\in\mathcal{T}_{\lambda}$ 且除了有限多个 $\lambda$ 外, $\mathcal{O}_{\lambda}=X_{\lambda}$.
\end{definition}

若所有 $X_{\lambda}$ 是同一个空间 $X$, 习惯上记 $\prod_{\lambda\in\Lambda}X_{\lambda}$ 为 $X^{\Lambda}$. 特别地, 若 $\mathbf{N}$ 表示自然数集, 则 $X^{\mathbf{N}}$ 是 $X$ 中的序列全体而 $\mathbf{R}^{X}$ 是以 $X$ 为定义域的实值函数全体. 若 $X$ 是度量空间而 $\Lambda$ 是可数的, 则 $X^{\Lambda}$ 上的乘积拓扑由度量诱导(见习题 16). 一般地, 若 $X$ 是度量空间而 $\Lambda$ 是不可数的, 乘积拓扑不能通过度量诱导. 例如, $\mathbf{R}^{\mathbf{R}}$ 上的乘积拓扑不能通过度量诱导(见习题 17). 我们将以下两个命题的证明留作练习.

\begin{proposition}\label{proposition:ch12-乘积逐点收敛}
令 $X$ 为拓扑空间. 序列 $\{f_{n}:\Lambda\to X\}$ 在乘积空间中收敛于 $f$ 当且仅当对每个 $\Lambda$ 中的 $\lambda$, $f_{n}(\lambda)$ 收敛于 $f(\lambda)$. 因此, 序列关于乘积拓扑的收敛性是逐点的收敛性.
\end{proposition}

\begin{proposition}\label{proposition:ch12-乘积弱拓扑}
拓扑空间的笛卡儿积 $\prod_{\lambda\in\Lambda}X_{\lambda}$ 上的乘积拓扑是与投影族 $\{\pi_{\lambda}:\prod_{\lambda\in\Lambda}X_{\lambda}\to X_{\lambda}\}_{\lambda\in\Lambda}$ 相联系的弱拓扑. 即它是笛卡儿积上使得所有投影映射连续的拓扑中具有最少集合个数的拓扑.
\end{proposition}

本节的中心部分是 Tychonoff 乘积定理, 它告诉我们紧拓扑空间的乘积 $\prod_{\lambda\in\Lambda}X_{\lambda}$ 是紧的. 该结果对指标空间 $\Lambda$ 没有限制. 在准备证明该定理前, 我们首先证明两个关于具有有限交性质的集族的引理.

\begin{lemma}\label{lemma:ch12-极大有限交}
令 $\mathcal{A}$ 为 $X$ 的具有有限交性质的子集族. 则存在 $X$ 的包含 $\mathcal{A}$ 的子集族 $\mathcal{B}$, 它具有有限交性质, 且关于该性质是最大的, 即不存在真包含 $\mathcal{B}$ 的具有有限交性质的 $X$ 的子集族.
\end{lemma}

\begin{proof}
考虑 $X$ 的包含 $\mathcal{A}$ 且具有有限交性质的子集族 $\mathcal{F}$. 用包含关系将 $\mathcal{F}$ 排序. $\mathcal{F}$ 的每个线性序子族 $\mathcal{F}_{0}$ 有由属于 $\mathcal{F}_{0}$ 中的任何族的集合组成的上界. 根据 Zorn 引理, 存在 $\mathcal{F}$ 的最大元. 该最大元是具有引理结论描述的性质的族.
\end{proof}

\begin{lemma}\label{lemma:ch12-极大有限交性质}
令 $\mathcal{B}$ 为 $X$ 的子集族. 它关于有限交性质是最大的. 则 $\mathcal{B}$ 中的有限个集合的交仍然在 $\mathcal{B}$ 中, 且 $X$ 的每个与 $\mathcal{B}$ 中的每个集有非空交的子集也属于 $\mathcal{B}$.
\end{lemma}

\begin{proof}
令 $\mathcal{B}'$ 为 $\mathcal{B}$ 中集合的有限交所成集合的全体. 则 $\mathcal{B}'$ 是包含 $\mathcal{B}$ 的具有有限交性质的族. 因此根据 $\mathcal{B}$ 关于包含关系的最大性, $\mathcal{B}'=\mathcal{B}$. 现在假定 $C$ 是 $X$ 的子集, 它与 $\mathcal{B}$ 的每个成员有非空交. 由于 $\mathcal{B}$ 包含 $\mathcal{B}$ 中集合的有限交, 得出 $\mathcal{B}\cup\{C\}$ 具有有限交性质. 根据 $\mathcal{B}$ 关于包含关系的最大性, $\mathcal{B}\cup\{C\}=\mathcal{B}$, 因此 $C$ 是 $\mathcal{B}$ 的成员.
\end{proof}

\begin{theorem}[Tychonoff 乘积定理]\label{theorem:ch12-Tychonoff}
令 $\{X_{\lambda}\}_{\lambda\in\Lambda}$ 为由集合 $\Lambda$ 标记的紧拓扑空间族. 则具有乘积拓扑的笛卡儿积 $\prod_{\lambda\in\Lambda}X_{\lambda}$ 也是紧的.
\end{theorem}

\begin{proof}
令 $\mathcal{F}$ 为 $X=\prod_{\lambda\in\Lambda}X_{\lambda}$ 的具有有限交性质的 $X$ 的闭子集族. 我们必须证明 $\mathcal{F}$ 有非空交. 由\cref{lemma:ch12-极大有限交}, 存在 $X$ 的子集族 $\mathcal{B}$(不必是闭的), 它包含 $\mathcal{F}$ 且关于有限交性质是最大的. 固定 $\lambda\in\Lambda$. 定义
\begin{align*}
\mathcal{B}_{\lambda}=\{\pi_{\lambda}(B)\mid B\in\mathcal{B}\}.
\end{align*}
则 $\mathcal{B}_{\lambda}$ 是 $X_{\lambda}$ 的具有有限交性质的子集族, $\mathcal{B}_{\lambda}$ 的成员的闭包族也如此. 根据 $X_{\lambda}$ 的紧性, 存在点 $x_{\lambda}\in X_{\lambda}$ 使得
\begin{align*}
x_{\lambda}\in\bigcap_{B\in\mathcal{B}}\overline{\pi_{\lambda}(B)}.
\end{align*}
定义 $x$ 为 $X$ 中的点, 其第 $\lambda$ 个坐标是 $x_{\lambda}$. 我们宣称:
\begin{align}
x\in\bigcap_{F\in\mathcal{F}}F.\label{align:ch12-Tychonoff交}
\end{align}
事实上, 点 $x$ 具有性质: 对每个指标 $\lambda$ 与每个 $B\in\mathcal{B}$, $x_{\lambda}$ 是 $\pi_{\lambda}(B)$ 的闭包点. 因此
\begin{align}
x\text{ 的每个次基本邻域 }\mathcal{N}_{x}\text{ 与 }\mathcal{B}\text{ 中的每个集合 }B\text{ 有非空交}.\label{align:ch12-次基本邻域}
\end{align}
从 $\mathcal{B}$ 的最大性与\cref{lemma:ch12-极大有限交性质}, 我们得出 $x$ 的每个次基本邻域属于 $\mathcal{B}$. 再次用\cref{lemma:ch12-极大有限交性质}, 我们得出 $x$ 的每个基本邻域属于 $\mathcal{B}$. 但 $\mathcal{B}$ 有有限交性质且包含族 $\mathcal{F}$. 令 $F$ 为 $\mathcal{F}$ 中的集合. 则 $x$ 的每个基本邻域与 $F$ 有非空交. 于是 $x$ 是闭集 $F$ 的闭包点, 因而 $x$ 属于 $F$. 因此\eqref{align:ch12-Tychonoff交}成立.
\end{proof}

\begin{exercise}\label{exercise:ch12-11}
证明具有乘积拓扑的任意 Tychonoff 空间族的乘积也是 Tychonoff 的.
\end{exercise}

\begin{exercise}\label{exercise:ch12-12}
证明具有乘积拓扑的任意 Hausdorff 空间族的乘积也是 Hausdorff 的.
\end{exercise}

\begin{exercise}\label{exercise:ch12-13}
考虑 $\mathbf{R}$ 的 $n$ 重笛卡儿积.
\begin{align*}
\mathbf{R}^{n}=\overbrace{\mathbf{R}\times\mathbf{R}\times\cdots\times\mathbf{R}}^{n}.
\end{align*}
证明该乘积拓扑与由欧氏度量诱导的 $\mathbf{R}^{n}$ 上的度量拓扑相同.
\end{exercise}

\begin{exercise}\label{exercise:ch12-14}
令 $(X,\rho_{1})$ 与 $(Y,\rho_{2})$ 为两个度量空间. 证明 $X\times Y$ 上的乘积拓扑与由乘积度量
\begin{align*}
\rho((x_{1},y_{1}),(x_{2},y_{2}))=\sqrt{[\rho_{1}(x_{1},x_{2})]^{2}+[\rho_{2}(y_{1},y_{2})]^{2}}
\end{align*}
诱导的拓扑相同, 其中 $X$ 和 $Y$ 具有由它们相应的度量诱导的拓扑.
\end{exercise}

\begin{exercise}\label{exercise:ch12-15}
证明: 若 $X$ 是具有度量 $\rho$ 的度量空间, 则
\begin{align*}
\rho^{*}(x,y)=\frac{\rho(x,y)}{1+\rho(x,y)}
\end{align*}
也是 $X$ 上的度量且它与 $\rho$ 诱导出相同的拓扑.
\end{exercise}

\begin{exercise}\label{exercise:ch12-16}
考虑度量空间的可数族 $\{(X_{n},\rho_{n})\}_{n=1}^{\infty}$. 对这些集的笛卡儿积 $X=\prod_{n=1}^{\infty}X_{n}$, 定义 $\rho:X\times X\to\mathbf{R}$ 为
\begin{align*}
\rho(x,y)=\sum_{n=1}^{\infty}\frac{\rho_{n}(x_{n},y_{n})}{2^{n}[1+\rho_{n}(x_{n},y_{n})]}.
\end{align*}
用前一个习题证明 $\rho$ 是 $\prod_{n=1}^{\infty}X_{n}$ 上的度量, 它诱导 $X$ 上的乘积拓扑, 其中每个 $X_{n}$ 具有由度量 $\rho_{n}$ 诱导的拓扑.
\end{exercise}

\begin{exercise}\label{exercise:ch12-17}
考虑具有乘积拓扑的集合 $X=\mathbf{R}^{\mathbf{R}}$. 令 $E$ 为 $X$ 在可数集上取值为 $0$ 而在其他集合上取值为 $1$ 的函数组成的子集. 令 $f_{0}$ 为恒等于零的函数. 则显然 $f_{0}$ 是 $E$ 的闭包点. 但 $E$ 中不存在收敛到 $f_{0}$ 的序列 $\{f_{n}\}$, 由于对 $E$ 中的任何序列 $\{f_{n}\}$, 存在某个 $x_{0}\in\mathbf{R}$ 使得对所有 $n$, $f_{n}(x_{0})=1$, 从而序列 $\{f_{n}(x_{0})\}$ 不收敛到 $f_{0}(x_{0})$. 这表明, $X=\mathbf{R}^{\mathbf{R}}$ 不是第一可数的且因此不是可度量化的.
\end{exercise}

\begin{exercise}\label{exercise:ch12-18}
令 $X$ 表示具有两个元素的离散拓扑空间. 证明 $X^{\mathbf{N}}$ 同胚于 Cantor 集.
\end{exercise}

\begin{exercise}\label{exercise:ch12-19}
用 Tychonoff 乘积定理与实数的闭有界区间的紧性, 证明 $\mathbf{R}^{n}$ 的任何闭有界子集是紧的.
\end{exercise}

\begin{exercise}\label{exercise:ch12-20}
给出以下断言的直接证明: 若 $X$ 是紧的而 $I$ 是闭有界区间, 则 $X\times I$ 是紧的. (提示: 令 $\mathcal{U}$ 为 $X\times I$ 的开覆盖, 且考虑 $t\in I$ 的最小值, 使得对每个 $t'<t$ 集合 $X\times[0,t']$ 可被 $\mathcal{U}$ 中的有限个集合覆盖. 用 $X$ 的紧性证明 $X\times[0,t]$ 也可被 $\mathcal{U}$ 中的有限个集合覆盖, 且若 $t<1$, 则对某个 $t'>t$, $X\times[0,t']$ 也可被 $\mathcal{U}$ 中的有限个集合覆盖.)
\end{exercise}

\begin{exercise}\label{exercise:ch12-21}
证明可数个序列紧拓扑空间的乘积是序列紧的.
\end{exercise}

\begin{exercise}\label{exercise:ch12-22}
单位区间的乘积 $I^{\Lambda}$ 称为(广义)方体. 证明每个紧 Hausdorff 空间 $X$ 同胚于某个方体的闭子集. (令 $\mathcal{F}$ 为 $X$ 上取值于 $[0,1]$ 的连续实值函数族. 令 $Q=\prod_{f\in\mathcal{F}}I_{f}$. 则由于 $X$ 是正规的, 从 $X$ 到 $Q$ 的将 $x$ 映到第 $f$ 个坐标是 $f(x)$ 的点的映射 $g$ 是一对一、连续且闭的象.)
\end{exercise}

\begin{exercise}\label{exercise:ch12-23}
令 $Q=I^{\Lambda}$ 为方体. 令 $f$ 为 $Q$ 上的连续实值函数. 则给定 $\varepsilon>0$, 存在 $Q$ 上的连续实值函数 $g$ 使得 $|f-g|<\varepsilon$ 且 $g$ 是仅有有限个坐标的函数. (提示: 用有限个长度为 $\varepsilon$ 的区间覆盖 $f$ 的值域且观察这些区间的原象.)
\end{exercise}

\end{document}